\documentclass[12pt]{jlreq}
\usepackage{amsmath, amssymb}

\newcommand{\C}{\mathbb{C}}
\newcommand{\R}{\mathbb{R}}
\newcommand{\Z}{\mathbb{Z}}
\newcommand{\N}{\mathbb{N}}
\newcommand{\F}{\mathbb{F}}
\newcommand{\Prob}{\mathbb{P}}
\newcommand{\Exp}{\mathbb{E}}
\newcommand{\dd}{\,d}
\newcommand{\Ker}{\operatorname{Ker}}
\newcommand{\Impart}{\operatorname{Im}}
\newcommand{\Hom}{\operatorname{Hom}}

\begin{document}

以下のような未知関数 \(v(t,a),\ Q(t)\) に関する微分方程式を領域 \(t\ge0,\ a\ge0\) で考える:
\[
  \frac{\partial v(t,a)}{\partial t}+\frac{\partial v(t,a)}{\partial a}
  =-\bigl(\mu(a)+f_1(P(t))+f_2(Q(t))\bigr)v(t,a),
\]
\[
  v(t,0)=\int_0^\infty \beta(a)v(t,a)\dd a,
\]
\[
  P(t)=\int_0^\infty v(t,a)\dd a,
\]
\[
  \frac{dQ(t)}{dt}=-bQ(t)+g(P(t))Q(t).
\]
ここで，\(b\) は正の定数，\(\beta(a),\ \mu(a)\) は恒等的にはゼロではない非負の有界連続関数で，\(\beta\) は区間 \([0,\infty)\) で可積分であるとする。\(f_1(x),\ f_2(x),\ g(x)\) はいずれも \(x\in[0,\infty)\) で定義された非負狭義単調増大な連続関数で，\(f_1(0)=f_2(0)=g(0)=0\) である。また関数 \(\ell(a)\) を
\[
  \ell(a)=\exp\left(-\int_0^a\mu(\sigma)\dd\sigma\right)
\]
と定義するとき，
\[
  \int_0^\infty \beta(a)\ell(a)\dd a>1
\]
であると仮定する。また初期条件は正であると仮定する:
\[
  v(0,a)>0,\qquad Q(0)>0.
\]
以下では上記のシステムが \(t>0\) で非負解をもち，すべての \(t>0\) で \(v(t,\cdot)\in L^1(0,\infty)\)，\(\lim_{a\to\infty}v(t,a)=0\) であり，さらに以下が成り立つことを仮定してよい。
\[
  \frac{\partial}{\partial t}\int_0^\infty v(t,a)\dd a
  =
  \int_0^\infty \frac{\partial v(t,a)}{\partial t}\dd a.
\]

\begin{enumerate}
  \item 任意の \(t>0\) に関して，\(P(t)>0\) であることを示せ。
  \item 関数 \(w(t,a)\) を
  \[
    w(t,a)=\frac{v(t,a)}{P(t)}
  \]
  と定義する。このとき \(w(t,a)\) が，パラメータとして \(\beta,\mu\) のみを含むような偏微分方程式と境界条件をみたすことを示せ。
  \item (2) で求めた偏微分方程式が，ただ一つの正の平衡解を持つことを示し，その平衡値をもとめよ。
  \item (3) で求めた平衡解を \(w^*(a)\) とする。\(v(t,a)=P(t)w^*(a)\) であるとき，\(P(t),Q(t)\) のみをすべき常微分方程式システムを求めよ。
  \item (4) で求めた常微分方程式システムが，正の平衡解を持つための必要十分条件を求めよ。
\end{enumerate}

\end{document}